\documentclass[
t, % align text inside frame to t=top, b=bottom, c=center
10pt, % 8pt, 9pt, 10pt, 11pt, 12pt, 14pt, 17pt, 20pt available as text font
aspectratio=1610, % select your aspect ratio 4:3=43, 16:9=169, 16:10=1610
ngerman,
english,
%handout,
]{beamer}
\usetheme{Juelich}

\usepackage{babel}
\usepackage[utf8]{inputenc}
\usepackage{textcomp}
\usepackage{wrapfig}

\title{Cognitive Computational Neuroscience}
\subtitle{Kriegeskorte \& Douglas (2018) Nature Neuroscience}
\author{Journal Club~~\vrule width0.3pt~~Julia Sprenger}
\institute[My Institute]{INM-6}
\date{\today}
\titlegraphic{\includegraphics[height=0.8\paperheight]{material/drawing}}

\begin{document}
% only use \maketitle to set your titlepage
\fzjset{title page=image}
\maketitle

\part{Introduction}
% \makepart

\begin{frame}[label=history]
	\frametitle{Historical Background}
	\onslide<1->{
	Cognitive psychology
	\begin{itemize}
	 \item study of mental processes such as 'attention, language use, memory, perception, problem solving, creativity, and thinking'\footnote{\url{https://en.wikipedia.org/wiki/Cognitive_psychology}}\newline
	\end{itemize}
	}
	
	\onslide<2->{
	Allen Newell (1973)
	\begin{itemize}
	  \item 'You can't play 20 questions with nature and win'
	  \item Hypothesis testing needs to be complemented by the construction of comprehensive task-performing computational models\newline
% 	  \item tutorial is not complete yet, will be updated shortly
	\end{itemize}
	}
	
	\onslide<3->{
	Richard Feynman (1988)
	\begin{itemize}
	 \item 'What I cannot create, I do not understand'\newline
	\end{itemize}
    }
    
    \onslide<4->{
    \textbf{Cognitive Science (1980)}\\
    introduction of task-performing computational models (symbolic cognitive architectures, neural networks based on behavioural data)
    }
%     \onslide<4->{
%     \textbf{New Cognitive Neuroscience}\\
%     will be based on task-performing computational models, that explain how cognition arises from neurobiologically plausible dynamic components
%     }
\end{frame}

\begin{frame}
    \frametitle{Historical Background}
    Cognitive Neuroscience
    \begin{itemize}
     \item relate cognitive theories to the (human) brain using functional brain imaging
     \item mapping of cognitive functions to brain regions using
     \begin{itemize}
        \item EEG (1875)
        \item MEG (1968)
        \item PET (1950s)
        \item fMRI (1990)
     \end{itemize}
    \end{itemize}

\end{frame}

\begin{frame}
    \frametitle{Modern Imaging Techniques}
    \centering
    \includegraphics[height=0.7\textheight]{figures/figure1}
    \begin{itemize}
     \item a brain map does not reveal the computational mechanism
     \item but constrains for theory
     \item data-driven analysis provides only limited insights
    \end{itemize}

\end{frame}

\begin{frame}
    \frametitle{Advances in Cognitive Science}
    \begin{itemize}
     \item face-selective regions in human
     \item spacial clustering of face neurons in non-human primates
    \end{itemize}
\end{frame}


\begin{frame}[label=Approaches]
	\frametitle{Different Approaches}
	\onslide<1->{
	Cognitive sciences
	\begin{itemize}
	 \item interdisciplinary, scientific study of the mind and its processes\footnote{\url{https://en.wikipedia.org/wiki/Cognitive_sciencey}}
	 \item how humans learn \& think
	 \newline
	\end{itemize}
	}
	
	\onslide<2->{
	Computational Neuroscience
	\begin{itemize}
	  \item how brains adapt and compute\newline
	\end{itemize}
	}
	
	\onslide<3->{
	Artificial Intelligence
	\begin{itemize}
	 \item how to generate intelligent behaviour
	\end{itemize}
    }
\end{frame}


\begin{frame}
    \frametitle{Disciplines}
    \centering
    \includegraphics[height=0.7\textheight]{figures/figure2}
\end{frame}

\begin{frame}
    \frametitle{Recent advances}
    Cognitive Science
    \begin{itemize}
     \item top-down approach
     \item Bayesian cognitive models (optimal combination of prior knowledge with sensory evidence)
     \item unified perspective on probabilistic empirical inference
    \end{itemize}
    Computational Neuroscience
    \begin{itemize}
     \item bottom-up approach
     \item mathematical models of elementary computational components and their implementation with biological neurons
    \end{itemize}
    Artificial Intelligence
    \begin{itemize}
     \item demonstrates how component functions can be combined to create intelligent behaviour
     \item machine learning, deep neural networks
    \end{itemize}
    
    \textbf{Overarching Challenge}
    build solid bridges between theory and experiment
\end{frame}

\begin{frame}
    \frametitle{From Experiment Toward Theory}
    \begin{itemize}
     \item Models of connectivity and dynamics
     \begin{itemize}
        \item correlation among response time series \textrightarrow 'functional connectivity'
        \item anatomical connectivity
        \item graph theoretical analysis of connectivity measures
        \item generative modeling of dynamics \textrightarrow effective connectivity
        \begin{itemize}
         \item Dynamic Causal Modeling
         \item Granger Causality
         \item Transfer Entropy
        \end{itemize}
        \item can be applied within individuals (different states) and across individuals (disorders)
     \end{itemize}
     \item Decoding models
     \begin{itemize}
      \item 'what information is present in each brain region' (eg. stimulus orientation or type, face identity, belief related decisions)
      \item does not specify the representational format
      \item linearly decodable \textrightarrow \textit{explicit} information
      \item advanced decoding: reconstruction
     \end{itemize}
     \item Representational models
     \begin{itemize}
      \item comprehensive predictions about the representational space
      \begin{itemize}
       \item encoding models
       \item pattern component models (PCM)
       \item representational similarity analysis (RSA)
      \end{itemize}
     \item often based on description of stimuli
     \end{itemize}
    \end{itemize}
    These models do not reveal the computational mechanisms of information processing underlying some cognitive function
\end{frame}

\begin{frame}
    \frametitle{The many meanings of model}
    \begin{itemize}
     \item Data-analysis models (statistical description of measured variables)
     \item box-and-arrow models (information processing)
     \item oracle model (relies on information without describing the extraction from input)
     \item brain-computational model (mimics brain information processing, eg sensory encoding)
     \item ...
    \end{itemize}
\end{frame}


\begin{frame}
    \frametitle{From Theory to Experiment}
\end{frame}

\begin{frame}
    \frametitle{The Space of Process Models}
    \centering
    \includegraphics[height=0.7\textheight]{figures/figure3}
\end{frame}

\begin{frame}
    \frametitle{From Theory to Experiment}
    \begin{itemize}
     \item Neural network models
     \begin{itemize}
      \item provide a common language for building task-performing models that meet combined criteria from all three disciplines
      \item can be constructed as biologically plausible systems (feed forward, recurrent, ...)
      \item supervised and unsupervised learning
      \item requires large parametric complexity to capture world knowledge for intelligent behaviour
      \item overfitting problem \textrightarrow evalutation in term of generalization performance
      \item example of visual pathways: visual hierarchy is also formed in deep neural networks
     \end{itemize}
     \item Cognitive models
     \begin{itemize}
      \item high level description of cognitive processes without biological details
      \item production systems
      \begin{itemize}
       \item sequence of cognitive actions based on 'if ... then ...' rules
      \end{itemize}
     \item reinforcement learning
     \begin{itemize}
      \item learning to maximize long-term cumulative reward through interaction with environment eg. by value functions or policies
      \item requires balance between exploitation and exploration
      \item model-free control: learning by trial and error
      \item model-based control: enables intelligent action in novel situations
      \item episodic control: storage of past experiences
     \end{itemize}
     \item Bayesian models
     \begin{itemize}
      \item provides optimal behaviour under given data and priors
      \item challenging to learn a generative model from sensory data
     \end{itemize}
     \end{itemize}
    \end{itemize}
\end{frame}

\begin{frame}
    \frametitle{Why do cognitive science, computational neuroscience and AI need each other?}
    \begin{itemize}
     \item Cognitive science
     \begin{itemize}
      \item needs computational neuroscience to discover algorithms of information processing
      \item needs brain data to provide constrains for complex models
      \item progresses in close interaction with AI
     \end{itemize}
    \item Computational Neuroscience
    \begin{itemize}
     \item needs cognitive science to challenge it to engage higher-level cognition
     \item needs machince learning / AI to provide theoretical and technological basis for modeling functions with biologically plausible dynamical components
    \end{itemize}
    \item AI
    \begin{itemize}
     \item needs cognitive science to guide the engineering of intelligence (eg. benchmarks for tasks)
     \item needs computational neuroscience for algorithmic inspiration
     \item main challenge: integration of computational and statistical efficiency
    \end{itemize}
    \end{itemize}
\end{frame}

\begin{frame}
    \frametitle{Looking Ahead}
    The brain seamlessly merges bottom-up disriminative and top-down generative computations in perceptual inference, and model-free and model-based control.
    \begin{itemize}
     \item bottom up and top down
     \begin{itemize}
      \item most important funding initiatives for bottom-up approach: Human Brain Project (synthesize neuroscience data in biologically detailed dynamic models), US Brain Initiative (measurement and manipulation of brain activity)
      \item better understanding in the context of a prior theory
     \end{itemize}
    \item Marr's levels (1982)
    \begin{itemize}
     \item computational theory
     \item representation and algorithm
     \item neurobiological implementation
    \end{itemize}
    \item \textrightarrow convergence of the three disciplines on algorithms and representations
    \item example: child seeing escalator
    \begin{itemize}
     \item neural networks: recognition of visual elements
     \item bayesian nonparametric models: concept formation from single experience
    \end{itemize}
    \item power discrepancy: efficient statistical and computational implementation in the brain based on 20 Watt
    \item \textrightarrow need for collaborations between labs with complementary expertise \& open science culture
    \end{itemize}
\end{frame}



\begin{frame}
    \frametitle{Interaction Among Sharable Components}
    \begin{wrapfigure}{r}{0.4\textwidth}
     \includegraphics[width=0.4\textwidth]{figures/figure4}
    \end{wrapfigure}
    Tasks
    \begin{itemize}
     \item provide controlled environment for behaviour
     \item OpenAI's Gym, Universe, DeepMind's lab
     \item interactions with virtual stimuli, natural environment as games, mass participation
    \end{itemize}
    Data
    \begin{itemize}
     \item behavioural data during task performance
     \item structural \& functional brain data
    \end{itemize}
    Models
    \begin{itemize}
     \item task-performing computational models
     \item initially only performing specific task, but must ultimately generalize across tasks 
    \end{itemize}
    Tests
    \begin{itemize}
     \item comparison between computational models and brain data (within a specific task)
     \item conceptual challenge of level of comparison
    \end{itemize}
    Challenges
    \begin{itemize}
     \item design shareable tasks and provide human behaioral dataset to set bar for AI model performance
     \item share brain activity data to constrain models and quantatively compare to models
    \end{itemize}



\end{frame}


\end{document}
